\subsection{Planification et organisation initiale}

Pour notre projet, nous avons commencé par élaborer un diagramme de Gantt prévisionnel. Cette première étape, basée sur nos estimations, nous a permis de quantifier plus facilement la charge de travail à fournir sur ce mois et demi de développement et d'identifier les jalons principaux. De plus, pour faciliter le travail en équipe et la communication, nous avons décidé de nous réunir tous les jours pour travailler ensemble en présentiel. Le distanciel n'a été appliqué que lorsque la situation l'obligeait.





\subsection{Planification détaillée et jalons du projet}

Pour structurer notre travail, nous avons donc établi un diagramme de Gantt prévisionnel couvrant la période du 11 mai au 25 juin 2026. Cependant, lors de la réalisation, ces délais n'ont pas pu être strictement respectés.
Tout d'abord, la phase d'initialisation et de conception (prévue du 11 au 18 mai) s'est prolongée d’une semaine du au contexte de notre projet. Nous avons dû élargir nos recherches et revoir le choix de nos composants pour nous adapter.
Ensuite, la phase de développement parallèle (prévue du 19 mai au 8 juin) a été plus longue que prévu. En effet, nous n'avons pas reçu nos composants à temps, ce qui a décalé le début de nos implémentations physiques et nous a obligés à adapter notre rythme de travail.
Par conséquent, le temps alloué à la phase de tests et à la finalisation de nos livrables s'est retrouvé beaucoup plus court qu'estimé initialement. Malgré ce délai réduit, nous avons réussi à nous réorganiser pour achever l'ensemble des rendus et valider le fonctionnement du système avant la journée des projets le 25 juin 2026



\subsection{Outils de collaboration}
Dans un premier temps, nous avons mis en place un espace Google Drive permettant à chaque membre du groupe de déposer son travail. Cependant, au vu de la nature technique du projet, nous nous sommes très vite dirigés vers GitHub Desktop. Cet environnement s'est avéré beaucoup plus adapté pour la gestion de notre projet et le versioning, ce qui nous a permis de gagner un temps précieux. 



\subsection{Phasage et répartition des rôles} La répartition des rôles s'est faite de manière pragmatique, en fonction des tâches à accomplir. Durant les deux premières semaines, nous avons travaillé ensemble sur la recherche des composants. 

Une fois que nous avons été d'accord sur le choix des composants et sur la direction que devait prendre le projet, nous avons partagé le travail en fonction des compétences de chacun. Nous avons toutefois veillé à ce que chaque membre du groupe puisse ressortir du projet en ayant acquis de nouvelles connaissances. La répartition s'est articulée ainsi : 

Pôle Électronique : Benoît et Jérémy ont pris en charge la conception électrique et le routage des cartes.

Pôle Informatique : Hugues et Maxime se sont concentrés sur le développement logiciel.

TITRE ????? : Clarisse et Naïm ont adopté un rôle plus polyvalent. Ils sont intervenus en renfort sur les différents pôles selon les besoins, en plus de prendre en charge de la modélisation 3D des boîtiers. Ils ont également géré la réalisation des tâches annexes essentielles au projet, telles que la rédaction du rapport d'éthique et la conception du poster. 
